% nostarch — a technical-book class written for No Starch Press's house
% style by Boris Veytsman, released independently under the LPPL. This
% starter uses the class's typographic structure (half-title, formal
% title page, copyright page, index) but is NOT an official No Starch
% Press product and carries no publisher branding, address, or ISBN —
% those fields in upstream's own sample are the actual publisher's real
% contact details and are not appropriate to ship as boilerplate.
%
% nocfonts: No Starch's real books use commercial fonts (New Baskerville
% etc.) that are not present in this image; nocfonts tells the class to
% fall back to free substitutes instead, so this compiles under plain
% pdflatex without those commercial packages.
%
% Class documentation: texdoc nostarch | https://ctan.org/pkg/nostarch
% Licence: LPPL-1.3 or later, (c) 2008-2023 Boris Veytsman.
\documentclass[nocfonts]{nostarch}

\usepackage{natbib}
\usepackage{nshyper}
\usepackage{amsmath}
\makeindex

\begin{document}

\frontmatter

\author{Marcus Webb}
\nostarchlocation{}

\title{\emph{Practical Compiler Construction}}

\subtitle{Building a Language, One Pass at a Time}

\makehalftitle

\maketitle

\begin{copyrightpage}
\textbf{\sffamily\MakeUppercase{\emph{Practical Compiler Construction}}}

This is a sample starter document generated from the \texttt{nostarch}
\LaTeX{} class. It is not an official No Starch Press publication.
\end{copyrightpage}

\begin{dedicationpage}
For everyone who has ever stared at a shift-reduce conflict.
\end{dedicationpage}

\tableofcontents

\chapter*{Foreword}
This book grew out of a compilers course taught over several years. It
assumes you can already write ordinary programs, but have never
written a program that reads programs.

\mainmatter

\chapter{Lexical Analysis}
\label{ch:lexer}

The first stage of any compiler turns a stream of characters into a
stream of \emph{tokens}\index{tokens}. A lexer\index{lexer} recognizes
patterns — identifiers, numbers, keywords — using regular expressions
compiled down to a finite automaton.

\section{Regular Expressions to Automata}

Every regular expression can be converted to a nondeterministic finite
automaton (NFA)\index{finite automaton!nondeterministic} via Thompson's
construction, and every NFA to an equivalent deterministic finite
automaton (DFA)\index{finite automaton!deterministic} via the subset
construction.

\textbf{Note.} Most real lexer generators skip the NFA step entirely
and construct the DFA directly from the regular expression's syntax
tree.

\chapter{Parsing}
\label{ch:parser}

A parser\index{parser} consumes the token stream and builds a parse
tree according to a grammar. This chapter covers recursive-descent
parsing for LL(1) grammars and shift-reduce parsing for LR(1)
grammars\index{grammar!LR(1)}.

\textbf{Tip.} Recursive-descent parsers are easier to write by hand; LR
parsers are easier to generate from a grammar specification and handle
a strictly larger class of grammars\index{grammar!LL(1)}.

\section{Handling Ambiguity}

The classic dangling-else problem\index{dangling else} is a textbook
example of grammar ambiguity: an \texttt{if}-\texttt{then}-\texttt{else}
grammar resolves it by convention, binding each \texttt{else} to the
nearest unmatched \texttt{if}, per~\citet{aho2006compilers}. A rigorous
treatment of parser generator theory appears
in~\citet{grune2008parsing}.

\backmatter

\bibliographystyle{plainnat}
\bibliography{references}

\printindex

\end{document}
